\chapter{Fondements des mod\`eles de langage}
\label{ch:foundations}

\begin{quote}
\emph{<< Un mod\`ele de langage est une distribution de probabilit\'e sur des s\'equences de tokens. Tout le reste --- chatbots, g\'en\'eration de code, r\'esum\'e --- d\'ecoule de cette unique id\'ee. >>}
\end{quote}

% --- hook ---
En 2017, huit chercheurs de Google publient un papier au titre presque provocateur --- \emph{Attention Is All You Need}. Personne, ce jour-là, n'en mesure les conséquences. Cinq ans plus tard, leur architecture --- le Transformer --- a engendré GPT, Claude, LLaMA, Mistral, Gemini, et a déplacé l'attention de tout le monde de la recherche fondamentale vers un usage industriel quotidien. Ce qui rend ces modèles puissants n'est pourtant pas leur architecture. C'est qu'ils savent faire une chose, et une seule : prédire le token suivant. Ce chapitre s'attarde sur cette idée, parce que c'est elle qui ordonne tout ce qui suit --- du chatbot médical à l'agent qui écrit du code.
\section{Qu'est-ce qu'un mod\`ele de langage ?}

Un mod\`ele de langage attribue une probabilit\'e \`a une s\'equence de tokens $w_1, w_2, \ldots, w_n$ :
\[
P(w_1, w_2, \ldots, w_n) = \prod_{i=1}^{n} P(w_i \mid w_1, \ldots, w_{i-1})
\]

Cette factorisation dit : pr\'edire chaque token \'etant donn\'e tout ce qui pr\'ec\`ede. Un mod\`ele qui fait cela correctement peut g\'en\'erer du texte coh\'erent, traduire des langues, r\'epondre \`a des questions et \'ecrire du code --- parce que toutes ces t\^aches se ram\`enent \`a << quel est le token suivant ? >>.

\begin{definition}[Mod\`ele de langage]
Une fonction $P_\theta$ param\'etr\'ee par $\theta$ qui associe \`a une s\'equence de tokens une probabilit\'e. L'entra\^inement ajuste $\theta$ pour maximiser la vraisemblance des textes observ\'es.
\end{definition}

\begin{aitip}
Vous pouvez interagir avec un mod\`ele de langage d\`es maintenant. Ouvrez un notebook Colab et ex\'ecutez l'exemple GPT-2 de la Section~\ref{sec:first-generation}. Le mod\`ele compl\'etera n'importe quel texte que vous lui donnerez.
\end{aitip}

% ============================================================
\section{Tokenisation}

Les mod\`eles de langage ne traitent ni des caract\`eres bruts, ni des mots entiers. Ils traitent des \emph{tokens} --- des unit\'es de sous-mots apprises sur les donn\'ees.

\subsection{Pourquoi des sous-mots ?}

\begin{itemize}
  \item Niveau caract\`ere : vocabulaire min\'eral (< 300), mais s\'equences extr\^emement longues.
  \item Niveau mot : vocabulaire \'enorme (100\,000+), et les mots inconnus provoquent des \'echecs.
  \item Niveau sous-mot : vocabulaire mod\'er\'e (32\,000--100\,000), g\`ere n'importe quel mot, garde les s\'equences courtes.
\end{itemize}

\subsection{Byte Pair Encoding (BPE)}

BPE part des caract\`eres individuels et fusionne it\'erativement la paire adjacente la plus fr\'equente :

\begin{enumerate}
  \item D\'epart : \texttt{["l", "o", "w", "e", "r"]}
  \item Paire la plus fr\'equente : \texttt{("l", "o")} $\rightarrow$ fusion en \texttt{"lo"}
  \item Continuer jusqu'\`a atteindre la taille de vocabulaire vis\'ee.
\end{enumerate}

\begin{minted}{python}
# Tokenisation avec tiktoken (tokeniseur GPT)
import tiktoken

enc = tiktoken.get_encoding("cl100k_base")  # tokeniseur GPT-4
text = "Generative AI transforms how we create content."
tokens = enc.encode(text)
print(f"Text: {text}")
print(f"Token IDs: {tokens}")
print(f"Tokens: {[enc.decode([t]) for t in tokens]}")
print(f"Number of tokens: {len(tokens)}")
\end{minted}

\subsection{Tokenisation WordPiece}

WordPiece (utilis\'e par BERT) ressemble \`a BPE mais s\'electionne les fusions selon un crit\`ere de vraisemblance. La biblioth\`eque \texttt{tokenizers} de HuggingFace prend en charge les deux :

\begin{minted}{python}
from transformers import AutoTokenizer

# Tokeniseur BPE (GPT-2)
gpt2_tok = AutoTokenizer.from_pretrained("gpt2")
print(gpt2_tok.tokenize("Generative AI is transforming research."))

# Tokeniseur WordPiece (BERT)
bert_tok = AutoTokenizer.from_pretrained("distilbert-base-uncased")
print(bert_tok.tokenize("Generative AI is transforming research."))
\end{minted}

\begin{exercise}
\begin{enumerate}
  \item Tokenisez la phrase << Pneumonoultramicroscopicsilicovolcanoconiosis is a lung disease >> avec les tokeniseurs GPT-2 et DistilBERT. Comparez le nombre de tokens et les d\'ecoupages en sous-mots.
  \item Avec \texttt{tiktoken}, tokenisez un paragraphe en anglais, en fran\c{c}ais et en chinois. Quelle langue produit le plus de tokens pour un contenu \'equivalent ? Pourquoi ?
  \item Comptez les tokens dans les 1000 premiers caract\`eres d'un article Wikip\'edia. Quelle est la relation entre nombre de tokens et nombre de caract\`eres ?
\end{enumerate}
\end{exercise}

% ============================================================
\section{Plongements (embeddings)}

Les tokens sont des entiers. Les r\'eseaux de neurones ont besoin de vecteurs continus. Un \emph{plongement} (embedding) associe \`a chaque identifiant de token un vecteur dense dans $\mathbb{R}^d$ :

\[
\text{Plongement} : \{0, 1, \ldots, V-1\} \rightarrow \mathbb{R}^d
\]

o\`u $V$ est la taille du vocabulaire et $d$ la dimension du plongement (typiquement 768 ou plus).

\begin{minted}{python}
import torch
from transformers import AutoModel, AutoTokenizer

model_name = "distilbert-base-uncased"
tokenizer = AutoTokenizer.from_pretrained(model_name)
model = AutoModel.from_pretrained(model_name)

text = "The patient has a fever."
inputs = tokenizer(text, return_tensors="pt")
with torch.no_grad():
    outputs = model(**inputs)

# outputs.last_hidden_state shape: (batch, seq_len, hidden_dim)
print(f"Shape: {outputs.last_hidden_state.shape}")
print(f"Embedding dim: {outputs.last_hidden_state.shape[-1]}")
\end{minted}

\begin{aitip}
Dans les mod\`eles modernes, les plongements ne sont pas de simples tables de consultation. Apr\`es passage par les couches Transformer, les plongements de sortie sont \emph{contextuels} --- le vecteur pour << bank >> diff\`ere selon que le contexte est financier ou riverain.
\end{aitip}

% ============================================================
\section{La fonction softmax}

La derni\`ere couche d'un mod\`ele de langage produit un vecteur de \emph{logits} $z \in \mathbb{R}^V$ (un score par token de vocabulaire). Le softmax convertit les logits en probabilit\'es :

\[
\text{softmax}(z_i) = \frac{e^{z_i}}{\sum_{j=1}^{V} e^{z_j}}
\]

\begin{minted}{python}
import torch
import torch.nn.functional as F

# Logits simul\'es pour un vocabulaire de 5 tokens
logits = torch.tensor([2.0, 1.0, 0.5, -1.0, 3.0])
probs = F.softmax(logits, dim=0)
print("Logits:", logits.tolist())
print("Probabilities:", [f"{p:.4f}" for p in probs.tolist()])
print("Sum:", f"{probs.sum().item():.4f}")  # toujours 1.0
\end{minted}

% ============================================================
\section{Perplexit\'e}
\label{sec:perplexity}

La perplexit\'e mesure \`a quel point un mod\`ele de langage pr\'edit bien un ensemble de test. Plus c'est bas, mieux c'est.

\[
\text{PPL} = \exp\left(-\frac{1}{N}\sum_{i=1}^{N} \log P(w_i \mid w_{<i})\right)
\]

Intuition : si la perplexit\'e est 50, le mod\`ele est << aussi incertain que s'il choisissait uniform\'ement parmi 50 tokens >> \`a chaque \'etape.

\begin{minted}{python}
import torch
from transformers import GPT2LMHeadModel, GPT2Tokenizer

model = GPT2LMHeadModel.from_pretrained("gpt2")
tokenizer = GPT2Tokenizer.from_pretrained("gpt2")
model.eval()

text = "The Transformer architecture revolutionized natural language processing."
inputs = tokenizer(text, return_tensors="pt")

with torch.no_grad():
    outputs = model(**inputs, labels=inputs["input_ids"])
    loss = outputs.loss  # perte de cross-entropy
    perplexity = torch.exp(loss)

print(f"Loss: {loss.item():.4f}")
print(f"Perplexity: {perplexity.item():.2f}")
\end{minted}

\begin{exercise}
\begin{enumerate}
  \item Calculez la perplexit\'e de GPT-2 sur trois phrases : une phrase anglaise grammaticale, une version brouill\'ee, et une phrase en fran\c{c}ais. Expliquez les diff\'erences.
  \item \'Ecrivez une fonction \texttt{compute\_perplexity(model, tokenizer, text)} qui renvoie la perplexit\'e pour n'importe quel texte d'entr\'ee.
\end{enumerate}
\end{exercise}

% ============================================================
\section{Votre premi\`ere g\'en\'eration de texte}
\label{sec:first-generation}

\begin{minted}{python}
from transformers import pipeline

generator = pipeline("text-generation", model="gpt2")
result = generator(
    "Artificial intelligence will",
    max_new_tokens=50,
    do_sample=True,
    temperature=0.8,
)
print(result[0]["generated_text"])
\end{minted}

\begin{ethicsbox}
Les mod\`eles de langage apprennent sur des textes du web, qui contiennent biais, d\'esinformation et contenus nocifs. GPT-2 peut produire du texte raciste, sexiste, ou factuellement faux. \textbf{Ne d\'eployez jamais un mod\`ele de base sans filtres de s\'ecurit\'e.} Nous traiterons \'evaluation et s\^uret\'e au Chapitre~\ref{ch:eval-safety}.
\end{ethicsbox}

% ============================================================
\section{R\'esum\'e du chapitre}

\begin{itemize}
  \item Un mod\`ele de langage pr\'edit la probabilit\'e du token suivant \'etant donn\'e les tokens pr\'ec\'edents.
  \item La tokenisation (BPE, WordPiece) convertit le texte en s\'equences enti\`eres de sous-mots.
  \item Les plongements associent aux identifiants de tokens des vecteurs continus ; dans un Transformer, les plongements sont contextuels.
  \item Le softmax convertit les logits en une distribution de probabilit\'e sur le vocabulaire.
  \item La perplexit\'e mesure la qualit\'e d'un mod\`ele : plus la perplexit\'e est basse, meilleures sont les pr\'edictions.
  \item HuggingFace \texttt{transformers} fournit tokeniseurs, mod\`eles, et pipelines de g\'en\'eration pr\^ets \`a l'emploi.
\end{itemize}
